% Guión de la presentación EHMbrAIn (castellano).
% Build:  latexmk -pdf -outdir=build guion-presentacion.tex   (desde paper/slides/)
%
% Acompaña a EHMbrAIn-slides.pdf. Las diapositivas están en inglés: los títulos
% se citan aquí en inglés para poder localizarlas, la narración va en castellano.
%
% NOTA (repositorio público): el autor permanece anónimo en el repositorio.

\documentclass[11pt,a4paper]{article}

\usepackage[a4paper,top=2.2cm,bottom=2.2cm,left=2.6cm,right=2.6cm]{geometry}
\usepackage[utf8]{inputenc}
\usepackage[T1]{fontenc}
% es-noshorthands: los shorthands de babel-spanish ("~, "<, etc.) rompen
% tikz/siunitx y dieron problemas en este proyecto. Desactivados a propósito.
\usepackage[spanish,es-noshorthands,es-noquoting]{babel}
\usepackage{lmodern}
\usepackage{microtype}
\usepackage{amsmath,amssymb}
\usepackage{booktabs}
\usepackage{array}
\usepackage{xcolor}
\usepackage{enumitem}
\usepackage{fancyhdr}
\usepackage[hidelinks]{hyperref}

\definecolor{inkc}{HTML}{212529}
\definecolor{bluec}{HTML}{4263EB}
\definecolor{redc}{HTML}{A61E4D}
\definecolor{grayc}{HTML}{868E96}
\color{inkc}

\setlength{\parindent}{0pt}
\setlength{\parskip}{0.55em}
\linespread{1.06}

\pagestyle{fancy}
\fancyhf{}
\fancyhead[L]{\footnotesize\color{grayc} EHMbrAIn --- guión de presentación}
\fancyhead[R]{\footnotesize\color{grayc}\thepage}
\renewcommand{\headrulewidth}{0.2pt}
\renewcommand{\headrule}{\hbox to\headwidth{\color{grayc}\leaders\hrule height \headrulewidth\hfill}}

% --- marcadores de diapositiva ------------------------------------------
% \diapo{nº pdf}{título en pantalla}
\newcommand{\diapo}[2]{%
  \vspace{0.9em}%
  \noindent{\color{bluec}\rule{\linewidth}{0.6pt}}\par\vspace{-0.35em}%
  \noindent{\bfseries\color{bluec} Diapositiva #1}\quad{\bfseries #2}\par
  \vspace{-0.15em}}

% narración: lo que se dice en voz alta
\newenvironment{decir}%
  {\begin{list}{}{\leftmargin=1.1em \rightmargin=0em \listparindent=0pt
       \itemindent=0pt \parsep=0.5em \topsep=0.4em}\item[]\ignorespaces}%
  {\end{list}}

% acotación escénica: no se dice
\newcommand{\escena}[1]{{\itshape\color{grayc}[#1]}}
% énfasis de cifra clave
\newcommand{\cifra}[1]{\textbf{#1}}

\title{\vspace{-1.5cm}\textbf{EHMbrAIn}\\[2mm]
  \large Guión de la presentación}
\author{}
\date{\small Acompaña a \texttt{EHMbrAIn-slides.pdf} (44 páginas, 31 diapositivas numeradas)}

\begin{document}
\maketitle
\thispagestyle{fancy}
\vspace{-1.2cm}

\section*{Cómo usar este guión}

Los números de diapositiva son \textbf{páginas del PDF}. El texto sangrado es lo
que se dice en voz alta; lo que aparece \escena{así} son acotaciones, no se
pronuncian.

Las diapositivas están en inglés y sus títulos se citan aquí en inglés para poder
localizarlas de un vistazo. Los términos técnicos que aparecen en pantalla se
mantienen en inglés la primera vez que salen, con la traducción al lado.

\begin{center}\small
\begin{tabular}{@{}llr@{}}
  \toprule
  Bloque & Diapositivas & Tiempo \\
  \midrule
  Apertura & 1--2 & 2\,min \\
  1. El problema & 3--6 & 4\,min \\
  2. La física, y por qué el diagnóstico es difícil & 7--11 & 6\,min \\
  3. El banco de pruebas & 12--16 & 5\,min \\
  4. Qué entrega el GPA clásico & 17--21 & 5\,min \\
  5. El concurso pre-registrado & 22--28 & 7\,min \\
  6. Más allá del muro & 29--35 & 7\,min \\
  7. Qué llevarse & 36--40 & 4\,min \\
  \midrule
  \textbf{Total} & & \textbf{40\,min} $+$ 10 de preguntas \\
  \bottomrule
\end{tabular}
\end{center}

\textbf{Si vas justo de tiempo}, recorta en este orden: diapositiva 20
(detectabilidad), 28 (sim-to-real), 33 (piso prognóstico), 35 (economía).
\textbf{No recortes nunca} la 10, la 11, la 25 ni la 31: el muro de
observabilidad es la espina dorsal del argumento y sin esas cuatro la conclusión
de compra no se sostiene.

% =========================================================================
\section*{Apertura --- diapositivas 1--2 \hfill \normalsize\color{grayc} 2 min}
% =========================================================================

\diapo{1}{Portada}
\begin{decir}
Este es un estudio sobre una comparación que la literatura lleva tiempo haciendo
mal. El aprendizaje automático se aplica ya de forma rutinaria a la monitorización
de salud de motores, y de forma igual de rutinaria se publica que gana a los
métodos clásicos. Lo que voy a defender es que casi ninguna de esas comparaciones
es justa; que cuando haces una que sí lo es, la respuesta se parte de una manera
interesante; y que esa partición te dice exactamente dónde hay que gastar el
dinero.
\end{decir}

\diapo{2}{Where this talk goes}
\begin{decir}
Voy a trabajar en dos niveles todo el rato. Cada diapositiva abre con la
afirmación en una línea y debajo enseña el mecanismo que la produce. Si solo
sigues las líneas de arriba te llevas el argumento completo; el resto está ahí
para quien quiera comprobarlo.
\end{decir}
\escena{No leas el índice en voz alta. Señálalo, una frase, y sigue.}

% =========================================================================
\section*{1. El problema --- diapositivas 3--6 \hfill \normalsize\color{grayc} 4 min}
% =========================================================================

\diapo{4}{An engine degrades where nobody can look}
\begin{decir}
Esta es la situación física. Un turbofán comercial se queda atornillado a un ala
durante años. En ese tiempo el polvo y la contaminación ensucian los álabes del
compresor, la sección caliente sufre fluencia y se oxida, las holguras de punta
de álabe se abren, los sellos se desgastan. Nada de eso se ve. Nadie abre el
motor para mirar.

Lo que la aerolínea recibe a cambio es esto: un puñado de números por vuelo.
Velocidades de eje, temperatura de gases de escape, consumo de combustible. Esa
es toda la ventana.

\escena{señala la figura} Y este es el agregado que un operador vigila de verdad:
el margen de temperatura de gases de escape, la distancia al límite certificado.
Los dientes de sierra son lavados de compresor. Cuando llega a cero el motor sale
del ala. Una visita a taller cuesta millones de dólares; una retirada no
programada le destroza la programación a una flota entera durante días.
\end{decir}

\diapo{5}{Two schools, and a broken comparison}
\begin{decir}
Hay dos familias de métodos que interpretan esa ventana.

La clásica viene de los años setenta: el \emph{gas path analysis}, el análisis del
camino de gases. Construyes un modelo termodinámico del motor y lo inviertes: a
partir de los síntomas, trabajas hacia atrás hasta saber qué componente se ha
degradado y cuánto. Alrededor de ese núcleo la industria construyó una caja de
herramientas madura: modelos de referencia, suavizado de tendencias, filtros de
Kalman, reglas expertas.

La segunda es el aprendizaje automático, sobre todo en la última década, sobre
todo en predicción de vida remanente, y sobre todo sobre bancos de pruebas
sintéticos.

\escena{baja el ritmo aquí: esta es la afirmación que motiva toda la charla} Y
aquí está el problema. A la IA casi nunca se la compara contra un pipeline
tradicional implementado con el mismo cuidado: mismos datos, mismas métricas,
mismo esfuerzo de ajuste. Lo que el campo tiene es un flujo continuo de
comparaciones contra hombres de paja. Lo cual significa que en realidad no
sabemos qué aporta la IA, bajo qué condiciones, ni a qué coste.
\end{decir}

\diapo{6}{What this project builds}
\begin{decir}
Así que este proyecto construye la comparación justa. Tres piezas.

Un gemelo termodinámico de un motor concreto --- el CFM56-7B, el que llevan la
mayoría de los 737NG --- construido únicamente a partir de datos públicos
certificados. Cien motores virtuales llevados hasta el fallo, en el formato de
datos exacto que recibe una aerolínea. Y las dos familias de métodos sobre datos
idénticos, métricas idénticas, particiones idénticas y presupuesto de ajuste
idéntico, juzgadas bajo un protocolo congelado antes de ver ningún resultado.

\escena{la idea clave de la charla: aterrízala} La razón de hacerlo sintético está
abajo en la diapositiva. En una flota real puedes preguntar «¿predijo bien el
método?». Lo que no puedes preguntar es «¿le echó la culpa al componente
correcto?», porque nadie conoce el estado interno verdadero de un motor montado en
un ala. Con verdad-terreno sintética sí puedes. Y la atribución es precisamente
donde el GPA clásico es más débil y donde la IA es menos transparente --- así que
es la pregunta que justifica construir un banco de pruebas entero.
\end{decir}

% =========================================================================
\section*{2. La física --- diapositivas 7--11 \hfill \normalsize\color{grayc} 6 min}
% =========================================================================

\escena{Este es el cimiento de bajo nivel. Si el público es aeronáutico, ve
rápido. Si es de aprendizaje automático, esta es la parte que más necesitan.}

\diapo{8}{The machine, in one picture}
\begin{decir}
Muy brevemente, la máquina. El aire entra por el fan. La mayor parte rodea el
núcleo y produce la mayor parte del empuje. El resto se comprime, se quema, y se
expande a través de dos turbinas que mueven los compresores.

El hecho estructural importante es que hay dos ejes mecánicamente independientes,
dos \emph{spools}. El de baja presión: fan y booster movidos por la turbina de
baja, girando a una velocidad que se llama N1. El de alta presión: compresor y
turbina de alta, a N2.

Los números de abajo son numeración de estaciones y van a reaparecer en
diapositivas posteriores. Y la frase del pie es la restricción bajo la que vive
toda esta charla: la aerolínea recibe cuatro canales. N1, N2, temperatura de
escape y consumo de combustible. Algunas flotas compran un paquete opcional con
presiones y temperaturas dentro del núcleo. La mayoría no.
\end{decir}

\diapo{9}{Gas path analysis: the model to invert}
\begin{decir}
Ahora la matemática, y es corta.

Describimos la condición interna del motor con dos números por turbomáquina. El
rendimiento isentrópico --- cuánto del trabajo ideal consigue de verdad; la
suciedad y el desgaste lo reducen. Y la capacidad de paso --- cuánto aire traga;
el ensuciamiento cierra un compresor, la erosión abre una turbina. Cinco
turbomáquinas, así que diez números. Ese vector es lo que mantenimiento quiere
saber, y no es medible.

Lo que sí es medible son las desviaciones de los canales medidos respecto de lo
que leería un motor sano en la misma condición de operación. Para las desviaciones
pequeñas del desgaste en servicio la relación es lineal: desviaciones medidas igual
a una matriz de influencia por el vector de salud, más sesgo de sensor, más ruido.

\escena{pausa} Y ahora mira la forma de ese problema. Diez incógnitas. Tres
medidas utilizables --- N1 no cuenta, es el parámetro de empuje comandado, no un
síntoma. Así que la matriz tiene rango tres. Hay infinitos estados de salud
compatibles con cualquier medida. La respuesta clásica regulariza: penaliza
desviaciones implausiblemente grandes y elige una solución. Eso funciona, y te
cuesta precisión de una forma que vamos a medir exactamente.
\end{decir}

\diapo{10}{Smearing: how a diagnosis flips to the wrong component}
\begin{decir}
Esto es lo que hace ese mal condicionamiento en la práctica, en un problema lo
bastante pequeño como para comprobarlo a mano.

Dos parámetros de salud, dos medidas. La verdad es una pérdida del uno por ciento
de rendimiento en el compresor de alta. Metes las medidas exactas en la inversión
y lo recuperas exactamente.

Ahora perturba una medida una décima de por ciento --- una desviación típica de
ruido de sensor. \escena{señala} La respuesta pasa a ser: compresor sano, turbina
degradada un uno por ciento. El diagnóstico ha saltado al componente equivocado
por completo.

El dibujo de la derecha explica por qué. Cada fallo produce su propio patrón de
movimiento de los indicadores; piénsalo como una flecha en el espacio de medidas.
Diagnosticar es preguntar a lo largo de qué flecha apunta la lectura observada.
Cuando dos flechas son casi paralelas, un empujón del tamaño del ruido mueve la
lectura de una a la otra. Esto se llama \emph{smearing}, embadurnamiento, y no hay
nada mal en el álgebra. Sencillamente, la información no está en las medidas.
\end{decir}

\diapo{11}{This engine's actual geometry}
\begin{decir}
Eso era un juguete. Este es el motor real.

Con el conjunto de sensores de cabina, el ángulo más pequeño entre dos firmas de
fallo en este motor es de \cifra{1,3 grados} --- entre el rendimiento del compresor
de alta y el rendimiento de la turbina de alta. Siete pares quedan por debajo de
quince grados, que es más o menos donde el ruido de sensor funde dos fallos en
uno. Rango tres para diez incógnitas.

Añade las cuatro sondas de estación y el rango sube a siete y el peor ángulo se
multiplica por cinco.

\escena{esta es la afirmación portante de la charla: dila despacio} Y repetimos
esto sobre una rejilla de seis puntos que cubre la envolvente de vuelo. La
estructura apenas se mueve: entre 1,2 y 1,4 grados en todas partes. O sea que la
ceguera no es una propiedad de una condición de vuelo concreta de la que puedas
escapar midiendo en otro sitio. Es intrínseca al conjunto de sensores. Quedaos con
eso, porque todo lo que viene en la segunda mitad de la charla es una consecuencia.
\end{decir}

% =========================================================================
\section*{3. El banco de pruebas --- diapositivas 12--16 \hfill \normalsize\color{grayc} 5 min}
% =========================================================================

\diapo{13}{The twin: certified data in, predictions out}
\begin{decir}
El gemelo está construido en pyCycle, la librería abierta de modelado de ciclo de
la NASA, transformando un ejemplo validado hacia este motor: como máximo tres
parámetros por paso, exigiendo convergencia en cada paso. Suena quisquilloso; es
la diferencia entre tener un modelo y pasar tres semanas depurando un solver de
Newton.

La calibración es la parte que merece atención. Usa exclusivamente datos medidos
certificados: la hoja de certificado de tipo, y la base de datos de emisiones de
la OACI, que publica algo silenciosamente valioso --- consumos de combustible
\textbf{medidos} en banco a nivel del mar, en cuatro regímenes de empuje, para
esta variante exacta de motor, gratis.

Y resolvemos los puntos de anclaje con el empuje impuesto. Así que el consumo del
modelo es una \textbf{predicción}, confrontada contra una medida a la que nunca se
ajustó. Ese es el examen honesto, y sale dentro del uno y medio al tres por ciento
en despegue, ascenso y aproximación. El consumo específico en crucero cae dentro
del dos por ciento del valor publicado sin haber sido nunca un objetivo.
\end{decir}

\diapo{14}{The fleet: six mechanisms, one hundred lives}
\begin{decir}
A partir de ese gemelo generamos cien motores viviendo vidas completas. Seis
mecanismos de degradación, cada uno una función en forma cerrada del ciclo de
vuelo, con magnitudes tomadas de la literatura de degradación de turbinas de gas.

\escena{recorre los paneles} Ensuciamiento saturando hacia una asíntota con
recuperaciones por lavado. Erosión lineal. Holgura de punta que se abre rápido y
luego despacio. Sección caliente acelerando al final --- rendimiento abajo,
capacidad de paso arriba: ese par es la firma de una tobera de turbina erosionada.
Daño por objeto extraño como escalones de Poisson. Y episodios de fallo agudo
discretos.

El detalle que importa después: la contribución de cada mecanismo se guarda por
separado. Así que la flota no registra solo cuánto está degradado un motor, sino
\textbf{por qué}.
\end{decir}

\diapo{15}{Making the benchmark hard on purpose}
\begin{decir}
Ahora bien: la enfermedad más común de un dataset sintético es ser
accidentalmente fácil, y un banco de pruebas que halaga a tu método no vale nada.
Así que pusimos un umbral por adelantado: un clasificador deliberadamente débil
--- una regresión logística simple sobre canales crudos --- tiene que
\textbf{fallar} en la tarea de aislamiento, por debajo del sesenta por ciento.

Primera pasada: sacó casi un ochenta y dos por ciento. El umbral falló.

\escena{este es un momento de credibilidad: no lo corras} Diagnosticamos por qué.
Nuestras etiquetas de fallo eran el mecanismo crónico dominante, y los mecanismos
crónicos co-evolucionan con la edad en todos los motores. Así que la etiqueta era
un proxy de la edad, y los niveles crudos de los canales codifican la edad
directamente. La tarea que habíamos definido no era aislamiento de fallos en
absoluto.

Había dos respuestas posibles: relajar el umbral, o arreglar el diseño.
Rediseñamos el dataset --- episodios de fallo discretos de un solo parámetro,
elegidos entre los pares confundibles --- y pasó con un cincuenta y cuatro por
ciento. Después, en una versión posterior de la flota, el umbral mordió
\textbf{otra vez} con un sesenta y dos por ciento, y escalamos las magnitudes de
fallo hacia abajo hasta que pasó.

Los dos fallos están en el informe con todo detalle, porque un umbral que no puede
fallar no certifica nada.
\end{decir}

\diapo{16}{Four safeguards, fixed before any contestant ran}
\begin{decir}
Cuatro salvaguardas, todas fijadas antes de que corriera ningún método.

Pre-registro: hipótesis, umbrales, particiones con hashes de fichero, tests
estadísticos, todo congelado bajo una etiqueta de git antes de la pasada
confirmatoria. Prestado de los ensayos clínicos: declaras qué te convencería antes
de mirar la respuesta.

Presupuesto de ajuste simétrico: cincuenta pruebas de búsqueda automática para
cada familia. Esta importa más de lo que parece. Cualquier asimetría de esfuerzo
convierte silenciosamente una comparación de métodos en una comparación de
esfuerzo, y esa es exactamente la enfermedad de la literatura.

Partición por motor, nunca por vuelo: los vuelos de un mismo motor son casi
duplicados, y partir por vuelo permite que un modelo memorice motores concretos.

Y una línea base de verdad fuerte: el pipeline tradicional construido con forma
industrial, cada regla experta justificada.

El bloque de abajo es la disciplina que hace creíble todo lo demás. Los números
existen en dos castas y nunca se mezclan. Exploratorios, producidos mientras
construyes. Confirmatorios, producidos una sola vez, después del congelado, sobre
datos tocados exactamente una vez. Solo los segundos sostienen un veredicto.
\end{decir}

% =========================================================================
\section*{4. Qué entrega el GPA clásico --- diapositivas 17--21 \hfill \normalsize\color{grayc} 5 min}
% =========================================================================

\escena{Esta sección es nueva respecto a cómo se suele presentar este trabajo, y
es el suelo cuantitativo de todo lo que viene después. No te saltes la 19.}

\diapo{18}{The experiment: honest forward, classical inverse}
\begin{decir}
Antes de comparar nada con nada, una pregunta que en esta literatura nadie parece
responder: ¿qué entrega de verdad el método de libro de texto sobre este motor?

El diseño importa. Para la dirección directa --- la que hace de motor --- usamos un
gemelo neuronal \textbf{no lineal} del modelo termodinámico, con toda la curvatura
de la física. Para la inversa --- la que hace de sistema de monitorización --- usamos
los mínimos cuadrados regularizados \textbf{lineales} sobre la matriz de
influencia. Dos instrumentos distintos a propósito: invertir la misma matriz que
generó los datos sería circular y exacto por construcción, y no nos diría nada.

Aquí va una ejecución. El compresor pierde un uno por ciento de rendimiento. La
cabina ve el consumo subir medio punto, la temperatura subir medio punto, la
velocidad del núcleo bajar una décima. Inviertes eso y recuperas \cifra{un cuarto}
del fallo. Repítelo quinientas veces con ruido de sensor fresco y el método nombra
el componente correcto el \cifra{veintidós por ciento} de las veces. Su respuesta
equivocada favorita es la turbina.
\end{decir}

\diapo{19}{What the inversion gives back}
\begin{decir}
Ahora sistemáticamente, un fallo cada vez, los dos conjuntos de sensores.

Panel izquierdo, la magnitud que vuelve. Con el conjunto de cabina, entre un dos y
un setenta y seis por ciento según el parámetro. Con las sondas de estación, entre
treinta y noventa.

El panel derecho es el que hay que mirar. El índice de \emph{smearing}: la parte de
la magnitud diagnosticada que aterriza sobre componentes que en realidad están
sanos. Mediana \cifra{0,85}. Seis séptimas partes de lo que el método reporta se
atribuye al sitio equivocado. Y fíjate en que las sondas de estación apenas lo
mejoran: más sensores te compran magnitud, no limpieza, porque el estimador sigue
repartiendo lo que no puede separar.

El caso extremo merece nombrarse. Capacidad de paso de la turbina de alta: la
cabina recupera \cifra{un uno coma nueve por ciento} --- el fallo es prácticamente
invisible --- y el conjunto extendido recupera un noventa. Eso es un factor
cuarenta y siete comprado con una sonda de presión a la salida del compresor.
\end{decir}

\diapo{20}{How small a fault can be named} \escena{recortable}
\begin{decir}
Traduciendo eso al número que pediría un operador: ¿cómo de pequeño puede ser un
fallo para que este paquete de sensores todavía lo nombre bien nueve de cada diez
veces?

Conjunto de cabina: lo consigue para \cifra{uno} de los diez parámetros. El
conjunto extendido lo consigue para siete, varios por debajo del uno por ciento.

Y los fallos de nuestra flota están entre el 0,3 y el 1,1 por ciento. O sea que la
mayoría de los fallos de este banco de pruebas caen por debajo del límite de
resolución del conjunto de cabina --- lo cual va a explicar, dentro de unos cinco
minutos, por qué el concurso de aislamiento acaba como acaba.
\end{decir}

\diapo{21}{A real engine, eight faults at once}
\begin{decir}
Los fallos individuales son el ejercicio de libro. Los motores reales se degradan
por todas partes a la vez, así que aquí hay un motor real de la flota al final de
su vida, diagnosticado.

En negro, la verdad. El motor está dominado por daño de sección caliente:
rendimiento de turbina cuatro puntos abajo, y la capacidad de paso abierta dos
coma cuatro.

El diagnóstico de cabina, en rojo, acierta el titular: rendimiento de turbina tres
coma tres abajo. Y luego falla justo en la física que lo confirmaría: la apertura
de capacidad vuelve como \cifra{cero coma seis} en vez de dos coma cuatro. Así que
la firma clásica de tobera erosionada sencillamente no es visible en el
diagnóstico. El conjunto extendido, en azul, sí la recupera.

\escena{la parte honesta: dila} Y los dos se inventan daño que no existe. La
pérdida real de rendimiento del fan es de menos de medio punto; ambos reportan el
triple. La capacidad de paso real de la turbina de baja es esencialmente cero;
ambos reportan alrededor de menos uno coma tres --- equivocados en el signo además
de en el tamaño. Una decisión de mantenimiento tomada sobre esos dos números abre
un módulo sano.
\end{decir}

% =========================================================================
\section*{5. El concurso pre-registrado --- diapositivas 22--28 \hfill \normalsize\color{grayc} 7 min}
% =========================================================================

\diapo{23}{Five hypotheses, one confirmatory pass}
\begin{decir}
Ahora sí, la comparación. Las dos familias ajustadas bajo el presupuesto congelado
de cincuenta pruebas, los veinte motores de test evaluados exactamente una vez,
Wilcoxon emparejado por motor, McNemar exacto para los resultados binarios,
corrección de Holm entre hipótesis.

Tres confirmadas, dos refutadas. \escena{deja que lo lean un momento} Y quiero
dejar claro que las dos refutaciones no son fallos del estudio: son dos de sus
cuatro resultados más útiles. Las voy a tomar en orden.
\end{decir}

\diapo{24}{H1 --- the story is the tuning, not the model}
\begin{decir}
Detección primero, y aquí la historia va de método, no de modelos.

Antes del ajuste, nuestro detector de IA sacaba un recall de 0,06 y escribimos en
el informe que aquello parecía un problema de diseño, no de umbral. El espacio de
búsqueda congelado contenía justamente la palanca que sospechábamos: la longitud
de las ventanas de características. Con cincuenta pruebas el optimizador estiró
esas ventanas, y el recall pasó de 0,06 a \cifra{0,48}, con un retardo mediano de
499 ciclos. Doce veces antes que el detector clásico ajustado, con el mismo número
de falsas alarmas.

De aquí salen dos cosas transferibles. La longitud de ventana es la primera perilla
de la que vive o muere cualquier detector de cambios. Y --- la incómoda --- la
configuración ajustada del lado tradicional se seleccionó sobre una partición de
validación con solo dos motores limpios, y generalizó peor que sus propios valores
por defecto sin ajustar. La selección de umbrales sobre muestras limpias pequeñas
es frágil para \textbf{cualquier} método. Presupuestad motores limpios en
consecuencia.
\end{decir}

\diapo{25}{H2 refuted --- the wall binds everyone}
\begin{decir}
Aislamiento. Esta es la hipótesis que más esperábamos confirmar, y queda refutada.

Sobre los trece episodios confundibles de test las dos familias empatan
exactamente. Treinta y uno por ciento cada una. Y el patrón de desacuerdo es
perfectamente simétrico: cuatro episodios que solo acierta la IA, cuatro que solo
acierta la regla clásica. Está muy cerca de ser una moneda al aire cuál de las dos
se lleva el punto.

\escena{esta es la conclusión central de la charla} Ahora bien, esto lo predijimos
antes de que corriera ningún modelo, a partir de la geometría de la matriz de
influencia: firmas separadas 1,3 grados en un espacio de medidas de rango tres.
Las dos familias están limitadas por la misma información, y ninguna cantidad de
modelado recupera información que los sensores nunca capturaron.

Lo cual convierte una pregunta de software en una pregunta de compras. Si el
aislamiento de fallos confundibles importa en tu operación, la solución es
instrumentación.
\end{decir}

\diapo{26}{H3 and H5 --- prognosis is where learning pays}
\begin{decir}
La prognosis es donde el aprendizaje se gana el sueldo, y de forma decisiva.

Error de vida remanente de 858 ciclos al noventa por ciento de la vida, contra
1981 de la extrapolación clásica ajustada. Tamaño de efecto 0,89: coge al azar el
error de un motor de cada método y el clásico es mayor alrededor del noventa y
cuatro por ciento de las veces.

Pero mira la forma más que el número. \escena{señala} Los errores clásicos son
anchos y están sesgados hacia el \textbf{optimismo}: predice más vida remanente de
la que hay, sistemáticamente, que es la dirección peligrosa. Los de la IA son
estrechos y casi insesgados. El mecanismo es simple: una recta no puede curvarse
para seguir la degradación acelerada de la sección caliente, y un modelo de
secuencia aprende esa curvatura de los datos.

Y la declaración de incertidumbre que viene con ello, que es la hipótesis cinco. El
intervalo conformal cubre el ochenta y ocho por ciento empíricamente contra un
noventa nominal, con más menos dos mil ciclos. La banda clásica sobre-cubre al
noventa y ocho por ciento pero mide más menos once mil quinientos. Esa es la
diferencia entre una ranura de taller planificable y un encogimiento de hombros.
\end{decir}

\diapo{27}{H4 refuted --- physics features are not free wins}
\begin{decir}
La segunda refutación, y la que más me gustaría que un profesional se llevara.

«Physics-informed» es una etiqueta popular. La implementación obvia es meterle a
la red tus estimaciones de estado basadas en física junto con los datos crudos.
Pre-registramos la predicción de que eso ayudaría sobre todo cuando los datos
escasean.

Empeoró las cosas \textbf{en todas las fracciones de datos}, y empeoró más
justamente donde predijimos que más ayudaría. El mecanismo se entiende a
posteriori: esas estimaciones de estado están fuertemente embadurnadas. Son diez
canales ruidosos y mutuamente correlacionados cuyo contenido de información ya
está presente en las tres desviaciones a partir de las cuales se calcularon. Con
siete motores de entrenamiento, diez dimensiones de entrada extra son varianza
pura.

Después reabrimos esto por una segunda vía independiente --- otra característica,
otra flota, física no lineal --- y volvió a quedar refutado. Dos negativos. La
inyección de física necesita un canal de información que los datos crudos no
tengan; concatenar salidas de un estimador no lo es.
\end{decir}

\diapo{28}{Does the ranking survive outside our own simulator?} \escena{recortable}
\begin{decir}
La objeción obvia a todo lo anterior es que nosotros construimos el mundo en el
que se está puntuando a estos métodos. Así que reejecutamos la pregunta de
prognosis sobre C-MAPSS, de la NASA, el banco de pruebas de referencia del campo,
en cuya generación no tuvimos ninguna mano.

Misma dirección, magnitud comparable: 14,9 ciclos contra 42,1. El orden no es un
artefacto de nuestro simulador.

Una divulgación, a la derecha: nuestro plan nombraba el banco sucesor, más grande,
cuya distribución de varios gigabytes rompe la regla de este proyecto de que todo
tiene que correr en un ordenador de escritorio. Sustituimos, y lo decimos en el
registro de decisiones en lugar de cambiar el plan en silencio.
\end{decir}

% =========================================================================
\section*{6. Más allá del muro --- diapositivas 29--35 \hfill \normalsize\color{grayc} 7 min}
% =========================================================================

\diapo{30}{Can observability be bought without hardware?}
\begin{decir}
El resultado de aislamiento terminaba en «compra sensores», que es una conclusión
cara. Así que buscamos salidas más baratas, y encontramos dos.

La matriz de influencia no es una sola matriz: cambia con la condición de vuelo.
Así que una ventana de instantáneas ordinarias tomadas en condiciones dispersas es
un conjunto de proyecciones distintas del mismo estado de salud, y fusionarlas es
un problema de tomografía.

La auditoría dio una respuesta de dos caras. \escena{señala la figura} La dispersión
gratis --- la variación de ambiente y potencia que te sale sin pagar --- restaura el
rango algebraico inmediatamente, pero los ángulos saturan casi de inmediato. La
variación gratuita es real pero delgada.

Entonces: salida uno. Tratar el \textbf{calendario de reportes rutinarios} como una
variable de diseño de observabilidad, lo cual, hasta donde sabemos, es un
planteamiento nuevo. Añadir un reporte periódico de crucero estabilizado a baja
potencia compra un setenta y nueve por ciento más de separabilidad en la
ambigüedad de la turbina. El coste es un boletín de procedimiento, no aviónica.

Salida dos. Fusiona la ventana con un modelo aprendido, pero conserva la inversión
física. Una red alimentada con las inversiones físicas vuelo a vuelo aísla al 0,65
contra el 0,17 del apilado clásico. La misma red alimentada con secuencias crudas
fracasa directamente al 0,12. La división del trabajo es la lección: deja que la
física haga la inversión que hace bien, y que la red haga la fusión temporal que
hace bien.

Y lo que ningún calendario arregla: el par rendimiento de compresor contra
rendimiento de turbina se queda por debajo de 1,8 grados en toda la envolvente.
Ese es fundamental.
\end{decir}

\diapo{31}{Only a real sensor breaks the wall}
\begin{decir}
Lo cual plantea la pregunta que te hará un proveedor: ¿y no podemos
\textbf{predecir} el sensor que falta a partir de los que ya tienes? Tenemos un
gemelo calibrado; que calcule una sonda de presión virtual.

Tres condiciones, pre-registradas. Solo cabina: 0,15. Cabina más canales de
estación \textbf{virtuales}: 0,15, con p igual a uno. Cabina más canales de estación
\textbf{reales}: \cifra{0,92}.

Esto es la desigualdad de procesamiento de datos hecha operación. Una cantidad
calculada a partir de los datos de cabina no puede llevar información que los datos
de cabina no tuvieran ya, así que no puede separar fallos que la cabina no puede
separar. Ningún proveedor de software, por sofisticado que sea, sustituye a la
medida que falta.

La nota práctica está abajo, y es una noticia inusualmente buena: los parámetros
que rompen el muro ya los mide el ordenador de control del motor. La barrera es
registrarlos y bajarlos a tierra, no instalar hardware.
\end{decir}

\diapo{32}{A certificate of what a diagnosis can know}
\begin{decir}
Esta es la parte que más me gustaría que sobreviviera a la charla.

Todo diagnóstico de camino de gases en la industria es una estimación puntual que
esconde lo que no puede saber. Entonces: acumula información de Fisher sobre las
condiciones que un motor individual voló de verdad, a través del gemelo calibrado.
Invierte. Obtienes la cota de Cramér-Rao --- la mejor precisión por dirección que
\textbf{cualquier} estimador insesgado puede alcanzar con los datos de ese motor.
Eso es un certificado por motor, resuelto por su historia, de lo que su diagnóstico
puede afirmar honestamente.

Análisis de observabilidad de turbinas de gas existen. Lo que nunca ha sido posible
es el panel de la derecha: \textbf{comprobar que el certificado es honesto}, porque
en un motor real el estado verdadero de los componentes es incognoscible.

Aquí sí es cognoscible. La precisión certificada ordena el error real del filtro
contra la verdad-terreno con una correlación de 0,70. Las direcciones que el
certificado llama identificables son las que el estimador acierta; las que llama
inobservables son donde la estimación no vale nada. Hasta donde sabemos, esa es la
primera garantía de identificabilidad validada contra verdad-terreno para
diagnóstico de camino de gases --- y es una validación que solo un banco de pruebas
con verdad sintética puede llevar a cabo.
\end{decir}

\diapo{33}{Is a large RUL error the method's fault, or the future's?} \escena{recortable}
\begin{decir}
El mismo truco, aplicado al futuro en lugar de al presente.

Antes de financiar un mejor prognóstico, un operador debería preguntarse: ¿mi error
es grande porque mi método es débil, o porque el futuro es genuinamente
incognoscible a partir de lo que puedo ver hoy? Esas dos cosas tienen remedios
opuestos, y confundirlas o desperdicia presupuesto o abandona un objetivo que
estaba al alcance.

Así que condicionamos sobre la salud \textbf{verdadera} actual de cada motor,
buscamos sus vecinos más próximos en el espacio de salud --- motores en una
condición presente casi idéntica --- y medimos la dispersión de sus vidas
remanentes verdaderas. Dos motores en el mismo estado que fallan con cientos de
ciclos de diferencia difieren por razones que ninguna medida presente contiene. Esa
dispersión es un suelo irreducible.

A mitad de vida, el \cifra{ochenta y siete por ciento} de la incertidumbre de vida
remanente es irreducible. La vida remanente, a primer orden, no está escrita en el
presente --- lo cual es una corrección sobria a bastante trabajo de prognosis. Pero
al final de la vida se abre una brecha reducible de un factor cuatro. Ahí es donde
un mejor prognóstico paga, y solo ahí.
\end{decir}

\diapo{34}{The determinability map: when does each approach win?}
\begin{decir}
Un eje más, porque todo lo anterior vive en un único ajuste de calidad de sensor.
Regeneramos la flota con ruido de sensor a la mitad, al uno, al doble y al
cuádruple --- misma semilla, mismos motores, mismos fallos, solo se mueve la calidad
de medida --- y volvimos a puntuar las dos familias ajustadas.

Tres hallazgos, de izquierda a derecha.

La prognosis es esencialmente inmune al ruido: el error de la IA es plano a lo
largo de un factor ocho en ruido mientras que la extrapolación clásica se degrada
un cuarenta por ciento. La ventaja \textbf{se ensancha} de 2,3 a 3,4 veces. O sea
que el caso de negocio no depende de la calidad de los sensores.

El panel central es el que cierra el argumento de esta charla. Reducir a la mitad
todas las sigmas de sensor --- una mejora cara --- deja el aislamiento confundible
en 0,31 y 0,38. Ni de lejos el 0,92 que entregan sensores \textbf{distintos}. Dos
firmas separadas 1,3 grados siguen separadas 1,3 grados por precisamente que las
midas. Así que la orden de compra es específica: sondas en estaciones nuevas, no
mejores sondas en las de siempre.

Y a la derecha, el detector clásico se cae por un precipicio --- recall de 0,48 a
cero --- mientras que la IA se degrada con elegancia. La robustez frente a la
calidad de instrumentación puede que sea la ventaja más transferible del
aprendizaje aquí.
\end{decir}

\diapo{35}{What it is worth --- and why it might be nothing} \escena{recortable}
\begin{decir}
Brevemente, cuánto vale, y este capítulo abre con sus propias reservas a propósito.

Monetizamos exactamente \textbf{un} mecanismo, el único que los resultados
sostienen: una prognosis menos sesgada que convierte retiradas no programadas en
programadas. Todo lo demás que un entusiasta podría contar se discute y se deja
deliberadamente fuera del total.

Para un operador de cien motores, la mediana son tres millones de dólares al año.
Pero hay que leer la dispersión, no la mediana: percentil diez, 0,9 millones;
percentil noventa, 8 millones; y negativo neto en un uno por ciento de los
sorteos.

Y la columna de la derecha es por qué podría no valer nada: el muro de aislamiento
sigue en pie, así que un operador que espere mejor diagnóstico de esta inversión se
llevaría un chasco; las falsas alarmas erosionan confianza y dinero; y la
comprobación sim-to-real preservó el orden pero no necesariamente la magnitud.
\end{decir}

% =========================================================================
\section*{7. Qué llevarse --- diapositivas 36--40 \hfill \normalsize\color{grayc} 4 min}
% =========================================================================

\diapo{37}{The split verdict}
\begin{decir}
Entonces, la respuesta a la pregunta con la que empezó la charla.

El aprendizaje gana donde la tarea es agregación estadística a lo largo del tiempo.
Prognosis --- y ese resultado es robusto entre tres arquitecturas y a lo largo de la
calidad de sensor. Instante de detección. Incertidumbre calibrada.

Nadie gana donde la tarea está limitada por lo que los sensores observaron. El
aislamiento confundible es un empate, y el arreglo son setenta y siete puntos
porcentuales que dan las sondas de estación reales, y exactamente cero de las
virtuales.

\escena{el reencuadre: esta es la idea de cierre} Y quiero tener cuidado con lo que
eso dice sobre la física, porque es fácil malinterpretarlo. La física no perdió.
Inyectar física dentro del modelo aprendido falló, dos veces. Pero la física
\textbf{diagnosticó los límites}: la geometría de la matriz de influencia predijo,
antes de que corriera ningún aprendizaje, exactamente dónde iba a fallar cada
método --- y acertó. Usar la física para saber qué es cognoscible funcionó
completamente.
\end{decir}

\diapo{38}{What an engineering organization should do on Monday}
\begin{decir}
Si os lleváis una sola diapositiva, llevaos esta. En orden de cuánto dinero mueve
cada punto.

Desplegad aprendizaje para prognosis primero: ese es el caso de negocio. Detección:
detección de cambios moderna con ventanas ajustadas; la victoria vino de buscar la
ventana, que es barato y auditable, no de la profundidad de la red. Aislamiento:
comprad sensores, no modelos --- y concretamente sondas en estaciones nuevas.
Rediseñad el calendario de reportes antes de comprar nada. Proteged el parque de
sensores: un termopar a la deriva corrompe a todos los consumidores aguas abajo,
sean aprendidos o clásicos. Desconfiad de comparaciones sin ajustar en cualquiera
de las dos direcciones. Y exigid pre-registro de las afirmaciones de vuestros
proveedores.

\escena{el último punto, con una pausa antes} Un veredicto partido es el aspecto
que tiene una evaluación honesta. Si un proveedor os enseña un pleno, eso debería
subir vuestra sospecha, no vuestra confianza.
\end{decir}

\diapo{39}{What transfers off this simulator --- and what does not}
\begin{decir}
Por último, la honestidad intelectual que este tipo de trabajo exige.

Lo que transfiere fuera de un simulador son órdenes, mecanismos y geometría: qué
método gana dónde y por qué, y la estructura de observabilidad que comparte
cualquier motor de esta clase. Todas las afirmaciones operativas que he hecho son
de ese tipo.

Lo que \textbf{no} transfiere son los números absolutos. Un error de 858 ciclos es
una propiedad de este mundo simulado y de sus ajustes de ruido. No estoy afirmando
que vuestra flota vaya a ver eso.

Y toda la base de evidencia se regenera en unos once minutos en un solo ordenador
de escritorio. Cada figura de estas diapositivas salió de un script; ninguna se
copió a mano.
\end{decir}

\diapo{40}{The split verdict is the result}
\begin{decir}
\escena{pausa. Deja la diapositiva dos segundos antes de hablar.}

El veredicto partido es el resultado. Una comparación de métodos que produce un
pleno normalmente está midiendo esfuerzo, no método.

Gracias. Encantado de responder preguntas.
\end{decir}

% =========================================================================
\section*{Preparación de preguntas}
% =========================================================================

Las objeciones probables, con respuesta honesta. No te marques un farol: todas
tienen respuesta real dentro del trabajo.

\vspace{0.4em}
\textbf{\color{redc}«Vuestra línea base tradicional es una recta. Eso es un hombre de paja.»}
\begin{decir}
Es justo, y lo comprobamos. La extrapolación lineal es lo que las aerolíneas usan
de verdad, así que es la línea base correcta para una afirmación de transferencia a
la industria; pero no es la frontera de la investigación. Así que implementamos un
prognóstico basado en similitud, el método clásico estándar en C-MAPSS, que sí
puede seguir una degradación no lineal. Reduce casi a la mitad el error clásico: de
1981 a 1118. La IA sigue ganando con 858. Así que la afirmación honesta tiene dos
escalas: 2,3 veces mejor que lo que los operadores tienen hoy desplegado, y 1,3
veces mejor que el mejor método clásico. Una ventaja del treinta por ciento, no un
orden de magnitud.
\end{decir}

\textbf{\color{redc}«Es todo sintético. ¿Por qué debería creerme nada?»}
\begin{decir}
Cuatro defensas, todas construidas antes de los experimentos titulares. El
generador está anclado a medidas certificadas a las que nunca se ajustó. El dataset
tuvo que pasar un umbral de dificultad que falló dos veces y forzó rediseños. La
conclusión central se volvió a probar sobre un banco de pruebas que no generamos
nosotros. Y decimos explícitamente que los números absolutos no transfieren: solo
transfieren órdenes, mecanismos y geometría.
\end{decir}

\textbf{\color{redc}«Veinte motores de test. Trece episodios confundibles. Eso no es nada.»}
\begin{decir}
Correcto, y es nuestra principal limitación estadística. Es la razón de que cada
test esté emparejado por motor, de que usemos tests exactos en lugar de asintóticos,
y de que reportemos tamaños de efecto e intervalos bootstrap en vez de solo
p-valores. Es también la razón de que una de nuestras propias hipótesis del eje de
ruido quedara refutada sobre un umbral más fino de lo que la muestra podía resolver
--- y lo reportamos en lugar de esconderlo.
\end{decir}

\textbf{\color{redc}«Habéis refutado vuestras propias hipótesis. ¿No es eso un proyecto fallido?»}
\begin{decir}
Al contrario. El pre-registro es lo que hace creíble una refutación, y las
refutaciones son los resultados que cambiaron nuestras recomendaciones. H2
convirtió una pregunta de software en una pregunta de compras. H4 mató una técnica
popular en esta tarea, dos veces. Un estudio que confirma todo lo que predijo suele
ser un estudio que movió los postes de la portería.
\end{decir}

\textbf{\color{redc}«¿Por qué no una arquitectura moderna grande?»}
\begin{decir}
Comprobamos que el resultado no es un artefacto de la arquitectura: una red
convolucional temporal y un Transformer baten a la línea base clásica esencialmente
por el mismo margen que el modelo recurrente. La elección de arquitectura apenas
importa aquí, y las redes más grandes no ayudaron: el techo es la información que
hay en las medidas, no el tamaño del modelo.
\end{decir}

\textbf{\color{redc}«¿Funcionaría esto en un motor del que no tenéis datos?»}
\begin{decir}
La receta sí. El gemelo se construye únicamente a partir de una hoja de certificado
de tipo y de una base de datos pública de emisiones, y las dos existen para todo
motor certificado. Eso es deliberado: el método es portable aunque nuestros números
concretos no lo sean.
\end{decir}

\textbf{\color{redc}«¿Cuál es la aportación más novedosa?»}
\begin{decir}
El certificado, en la diapositiva 32. Análisis de observabilidad hay; una garantía
de identificabilidad por motor y resuelta por su historia que además esté
\textbf{demostrada honesta contra el estado verdadero de los componentes} no la
hay, porque en hardware real esa comprobación es imposible. Es la única aportación
que necesitaba que existiera todo este banco de pruebas.
\end{decir}

\end{document}
